The calculation of the projectile trajectory is being done in MATLAB \cite{matlab2024b}. To do that all the target coordinates in the world frame, the radius of the projectile and the mass of the projectile were written to a text file which then is read by the MATLAB \cite{matlab2024b} script. The radius and mass of the projectile were defined in the main program and not hardcoded into the MATLAB \cite{matlab2024b} script to easier allow for the implementation of logic to choose between different projectiles later. The drag coefficient for a sphere \cite{dragCoefficient}, the air density \cite{airDensity} and gravitational acceleration vector \cite{gravity} were directly defined in MATLAB \cite{matlab2024b}. The release position vector was defined dynamically, depending on the target coordinates. This was done to find suitable inverse kinematic solutions within the specified joint limits for a wider range of target points (see table 2,3 section kinematic solutions). To simulate the trajectory between the release point and the target point an ordinary differential equation for the flight curve of the projectile was solved numerically with “ode45” \cite{ode45} for a specified starting velocity vector and flight time \cite{dragEquation}. The error for each solution was calculated by comparing the actual landing position for the solution with the defined target position and finding the distance between the two. To find the best possible solution the optimizer function “fmincon” \cite{fmincon} was used to find the starting velocity vector and flight time to minimize the landing error starting from an initial guess. The optimizer also allows for the specifications of constraints. These were chosen as minimum -1 m/s and maximum +1 m/s for the velocity in all axis directions and maximum 1 m/s for the overall magnitude to comply with the URs technical specifications \cite{ur5TechnicalSpecifications} and to avoid unrealistic velocities. The flight time minimum was set as 0.001 seconds for the flight time as this needs to be over 0 s \cite{fmincon}.
